\subsection{Art Thou He That Art to Come?}\label{sec:second-week}

On the Second Sunday the station is the Roman church called Holy Cross in
Jerusalem. The Introit begins \latin{Populus Sion}, ``People of Sion'', and so
the Mass opens by calling the Church by the name of God's people in the holy
city. Last Sunday the Church asked God to stir up his own power and come. Now
she asks him to stir up our hearts:

\begin{studyframe}
\latin{Excita, Domine, corda nostra ad praeparandas Unigeniti tui vias: ut, per
ejus adventum, purificatis tibi mentibus servire mereamur: Qui tecum vivit
et regnat in unitate.}
\medskip

Stir up our hearts, O Lord, to prepare the ways of Thine only-begotten Son,
that through His coming we may be worthy to serve Thee with purified
minds. Who livest.\par
\end{studyframe}

The Gospel puts a question in the mouth of the greatest of the prophets.

\begin{quote}
Now when John had heard in prison the works of Christ: sending two of his
disciples he said to him: Art thou he that art to come, or look we for
another?\par
{\smallskip\noindent\small ---Matthew 11:2--3\par}
\end{quote}

John is in prison, and he sends his question to Christ. Christ does not reply
with a title. He sends back an account of what is happening:

\begin{quote}
The blind see, the lame walk, the lepers are cleansed, the deaf hear, the dead
rise again, the poor have the gospel preached to them.\par
{\smallskip\noindent\small ---Matthew 11:5\par}
\end{quote}

The list echoes Isaiah. ``God himself will come and will save you'', the
prophet had said. ``Then shall the eyes of the blind be opened, and the ears of
the deaf shall be unstopped.'' Christ tells John: look at what is being done,
and you will know who has come. Then, after the dead are raised, comes one
sign more, and it is not a miracle at all. The poor have the gospel preached to
them.

He adds a warning, ``And blessed is he that shall not be scandalized in me.''
St Gregory the Great, in the homily the Church reads at Matins, sees where the
warning points:

\begin{quote}
The sight of so many signs and so many mighty works should have been a source
of wonder, and not a stumbling-block. And yet the unfaithful found these very
works a rock of offence, when they afterwards saw Him Who had worked so many
miracles dying on the Cross. Hence Paul saith: ``We preach Christ crucified, unto
the Jews a stumbling-block and unto the Gentiles foolishness.'' \dots\ For the
more humbling God hath undergone for man's sake, the more worthy is He that man
should worship Him.\par
{\smallskip\noindent\small ---St Gregory the Great, Homily 6 on the Gospels\par}
\end{quote}

The station church bears the name of that Cross. In the second week of Advent,
with Christmas still far off, the Church is already asking whether her children
will follow the Messiah when the miracles end at the Cross. The One who opens
the eyes of the blind will be bound and will die on a cross, and Gregory says
that this is the very reason to worship him.

When John's disciples have gone, Christ turns to the crowd and praises the man
in prison: ``What went you out into the desert to see? a reed shaken with the
wind?'' Gregory says the Lord asks it in order to deny it. A reed bends at every
breath of air, as the carnal mind leans to whichever side praise or blame blows
it; John was no reed, for no change of circumstance bent him from his
uprightness.

Christ answered John's question out of Isaiah, and St Paul in the Epistle of the
day says why the Scriptures were given: ``what things soever were written were
written for our learning: that, through patience and the comfort of the
scriptures, we might have hope.'' John in his prison had need of both. The
Epistle ends with Isaiah's root of Jesse, ``in him the Gentiles shall hope'', and
with a prayer: ``Now the God of hope fill you with all joy and peace in
believing''.

Christ's own sign of his coming is the good news brought to the poor, and the
alms of the season find their reason here. Give them to one person, by your
own hand: a neighbour who is alone, a sick man no one visits, a prisoner, as
John was a prisoner.

\subsubsection*{The Immaculate Conception}

On 8 December, in the first week of Advent or the second, the Church keeps the
feast of the Immaculate Conception of the Virgin Mary. The season waits for the
Son; the feast looks at the dwelling God made ready for him. Its Collect says
that God kept her from every stain by the death of her Son, foreseen: the Cross
whose name the station of the Second Sunday bears reached back to her before she
bore him. And it asks that we too, by her prayers, may come to God pure in
heart:

\begin{studyframe}
\latin{Deus, qui per immaculatam Virginis Conceptionem dignum Filio tuo
habitaculum praeparasti: quaesumus; ut, qui ex morte eiusdem Filii tui praevisa
eam ab omni labe praeservasti, nos quoque mundos eius intercessione ad te
pervenire concedas. Per eundem Dominum.}
\medskip

O God, Who, by the Immaculate Conception of the Virgin, didst prepare a worthy
dwelling for Thy Son, we beseech Thee that Thou, Who, by the death, foreseen by
Thee, of the same Thy Son, didst preserve her from all stain, wilt grant us
also, by her intercession, to come to Thee pure in heart. Through the same.\par
\end{studyframe}
